\section{Alert Email Flooding}
\label{sec:threat-spam}

The alert service sends an email whenever a sensor reading exceeds a configured
threshold. A sustained flood of fabricated sensor payloads can exhaust the alert
channel, fill operator mailboxes, and cause genuine critical events to go unnoticed.
An attacker in possession of the
\texttt{INTERNAL\_TOKEN} (obtained in the APT phase described in
Section~\ref{sec:threat-apt}) can submit fabricated sensor payloads that
continuously exceed all thresholds, forcing the alert service to generate email
at the maximum rate allowed by the application. The attack is demonstrated by
\texttt{scripts/demo\_spam.sh}, which requires the token extracted by the
APT demo to be present.

\subsection{Detection}

Detection is implemented with two Wazuh rules on the alert-service access logs.
Rule 100300 fires on each \texttt{email\_sent} event, producing a marker per
outbound alert. Rule 100301 fires when 10 or more such markers accumulate within
60 seconds, confirming an abnormal burst rate.

\begin{lstlisting}[style=xmlstyle, caption={Wazuh rules for alert email flood detection.}]
<rule id="100300" level="10">
  <if_sid>86600</if_sid>
  <field name="event_type">email_sent</field>
  <description>Alert email sent from
    alert-service</description>
  <group>spam,email,</group>
</rule>

<rule id="100301" level="14" frequency="10" timeframe="60">
  <if_matched_sid>100300</if_matched_sid>
  <description>Email flood: 10+ alert emails in
    60 seconds</description>
  <group>spam,attack,flood,</group>
  <options>alert_by_email</options>
</rule>
\end{lstlisting}

Unlike the DoS and APT rules, rule 100301 uses no \texttt{same\_field} or
\texttt{different\_field} constraint, because all alert emails originate from
the alert service itself, not from an external IP. The flood is an application-level
pattern (volume over time) rather than a per-source-IP pattern.

\subsection{Response}

Two independent mechanisms bound the flood, operating at different layers.

At the application layer, a circuit breaker in the alert service stops email
delivery after 12 consecutive alerts and enforces a 5-minute cooldown before
resuming. This limits the total number of emails delivered in any single flood
to approximately 12 before the channel is closed.

At the transport layer, Postfix enforces a hard rate limit of 30 messages per
minute, regardless of application-layer state. This provides an independent
backstop that prevents transport saturation even if the circuit breaker is
bypassed or misconfigured.

Rule 100301 does not trigger an IP block. The alert service is a legitimate
internal component. Blocking its IP would suppress all future alerting. The rule
provides observability and email notification at level 14, ensuring operators are
aware of the flood even as the circuit breaker is engaged.

\subsection{Evaluation}

Both response mechanisms function correctly. In testing, the circuit breaker
engages after 12 emails and the 5-minute cooldown prevents any further delivery.
The Grafana \emph{Spam Alerts} stat (rule 100301 counter) and \emph{Alert Emails
Sent} timeseries confirm the burst pattern and its cutoff.

The primary limitation is that legitimate alert bursts (arising from a real
multi-zone environmental incident) are indistinguishable from an application-layer
flood by rate alone. A threshold of 10 emails per 60 seconds may generate false
positives during genuine disaster scenarios. Additionally, the circuit breaker
threshold (12 emails) is a fixed configuration. Appropriate values depend on the
operational context and would require tuning in a production deployment.
